\begin{proof}
For $r\ge2s_*$ choose a retained dyadic parameter $s$ with $r/2<s\le r$. Each rectangular partition cell has both side lengths at most $r$ and meets at most four target cells. Given its occupancy, its contribution to the target pattern can be recorded by one of at most fifteen nonempty masks. Therefore
\[
 H(Z_r^{\mathrm{grid}})\le H(Z_s)+\log_2(15)\E N_s
 \le C m(R)/m(r),
\]
using fixed-factor regularity. The bounded remaining range follows by data processing from $H(A)\le Cm(R)$.
\end{proof}

\section{Uniform inputs for the percolation models}\label{sec:inputs}

We use the following two regimes. In both, rectangles have horizontal and vertical sides in the displayed Euclidean embedding and aspect ratio at most $\kappa$.
\begin{enumerate}[label=(\roman*)]
\item Critical bond percolation on $\mathbb Z^2$, with parameter $1/2$.
\item Site percolation on
\[
 \mathbb T=\{(i+j/2,\sqrt3j/2):i,j\in\mathbb Z\},
\]
with independent parameters $p_x\in[p_*,1-p_*]$, where $p_*\in[p_0,1/2]$, $p_0>0$ is fixed, and $R\le L_\varepsilon(p_*)$. Here $L_\varepsilon$ is the characteristic length defined in \cite[Section~3.1]{Nolin}, with fixed $\varepsilon\in(0,1/2)$ and $L_\varepsilon(1/2)=\infty$.
\end{enumerate}
The second regime includes critical triangular percolation and spatially varying near-critical parameters. Constants are uniform over the indicated parameter profiles, rectangle locations, and retained scales; they may depend on $p_0,\varepsilon,\kappa$. All statements concern $R\ge R_0$ for a fixed sufficiently large lattice cutoff. Numbered citations to Nolin~\cite{Nolin} refer to the published version.

To pass from lattice-coordinate RSW to the Euclidean rectangles used here, cover each fixed-aspect-ratio corridor by a bounded number of overlapping lattice-coordinate parallelograms of side at most $cR$, where $c>0$ is fixed and sufficiently small. Crossing and joining these parallelograms gives the required corridor crossings by RSW and positive association. Thus each near-critical RSW input is applied at a lattice scale below $L_\varepsilon(p_*)$, even after the change of coordinates.

Critical RSW, arm separation, quasi-multiplicativity, and \eqref{eq:armgap} are available in regime (i) from \cite[Section~2.2]{GPS}. The half-plane three-arm bound \eqref{eq:halfarm} holds with $\beta=2$ by \cite[equation~(4.12)]{GPS}; lattice rotations give all four side directions. For regime (ii), Nolin's sandwich-product formulation, arm separation, and arm comparison \cite[Definition~5, Theorems~11 and~27, Proposition~17]{Nolin} give uniform RSW and \eqref{eq:armactual}. The next lemma gives a half-plane estimate uniform in the boundary direction in both regimes.

\begin{lemma}[Independent central inversion]\label{lem:reflection}
In either regime, let $L=\infty$ for critical square-lattice bonds and
$L=L_\varepsilon(p_*)$ for the triangular-lattice sandwich product measures.
For every Euclidean half-plane $H$, every $z\in\partial H$, either alternating
three-arm color order, and $r_*\le r<t\le L$,
\begin{equation}\label{eq:reflection}
 \Prob\bigl(A_3^H(z;r,t)\bigr)
 \le C(r/t)^{1+\eta_6/2},
\end{equation}
where $\eta_6>0$ is a critical alternating six-arm gap for the model:
$\alpha_{6,\mathrm{crit}}(a,b)\le C(a/b)^{2+\eta_6}$.
The constants are uniform in the direction and position of the boundary,
including its microscopic phase.  Any fixed microscopic padding of the
half-plane or of the annulus is allowed, with constants depending on that
padding.  Bond arms use open primal paths and closed dual paths.
\end{lemma}
\begin{proof}
Fix one of the two orders and denote its event by $E$.
Write $H=\{x:n\cdot(x-z)\ge0\}$, where $n$ is its inward unit normal.
Identify a random coordinate with its site, or with its primal-edge midpoint.
Choose a fixed $d_0$ so that every coordinate used by $E$ belongs to
\[
 U=\{x:n\cdot(x-z)\ge-d_0\}.
\]
Here and below a geometric set also denotes the random coordinates it contains.
Let $\ell$ exceed the distance from a coordinate to any point of its site-path
or primal/dual-edge realization, including all fixed endpoint and padding
conventions.  Choose a central inversion $g(x)=2c-x$ satisfying
\begin{equation*}
 |c-z|\le D_0,\qquad
 n\cdot(c-z)<-d_0-2\ell-2,
\end{equation*}
where $D_0$ is fixed independently of $n,z$.

For triangular sites, take $c$ in $\mathbb T$.  The lattice has bounded
covering radius, so such a $c$ exists by choosing it close to a point a fixed
distance from $z$ in direction $-n$.  Then $g$ is an involutive lattice
isometry, and the coordinate map is $\tau(x)=g(x)$.

For square bonds, take
\[
 c\in(1/4,1/4)+(1/2)\mathbb Z^2,
 \qquad 2c=a+(1/2,1/2),\quad a\in\mathbb Z^2.
\]
This set of centers also has bounded covering radius.  The map
$g(x)=a+(1/2,1/2)-x$ interchanges the primal lattice and its dual.  For a
primal edge $e$, define
\[
 \tau(e)=g(e)^*,
\]
where the star denotes the perpendicular primal edge crossing the dual edge
$g(e)$.  Taking a crossing edge commutes with $g$, so $\tau^2(e)=e$.
Thus $\tau$ is a bijection of the random bond coordinates, with midpoint
$g(\operatorname{mid}(e))$.  In particular, it exchanges horizontal and
vertical coordinates.  Moreover,
$g(e^*)=\tau(e)$: a primal path is carried to a dual path, and a dual path
to a primal path.

In both models the coordinate sets $U$ and $\tau(U)$ are disjoint.  Indeed,
their geometric half-planes are separated by more than $4\ell+4$.
The same separation, after allowing distance $\ell$ from the coordinates,
keeps all realizations of the two groups of paths disjoint.
Define a product measure $\nu$ by retaining $p_x$ on $U$, setting the
parameter at $\tau(x)$ to $1-p_x$ for every $x\in U$, and using $1/2$
on the remaining coordinates.  This remains in the triangular sandwich
class; in the bond model it is just the critical fair product measure.

Let $E'$ be the $\tau$-image of $E$ with every coordinate state reversed.
It is supported on $\tau(U)$ and has $\nu(E')=\Prob(E)$.
For bonds the state reversal gives exactly the required geometric colors:
an open primal path becomes a closed dual path, and a closed dual path
becomes an open primal path.  Independence now gives
\begin{equation}\label{eq:reflectsquare}
 \Prob(E)^2=\nu(E\cap E').
\end{equation}

On $E\cap E'$, truncate the six paths to a common annulus. Choose a lattice
center $q$ within bounded distance of $z$; both $z$ and $g(z)$ are then
within a fixed distance $D$ of $q$. Let $S_a(q)$ denote the balls in the
lattice norm used for the arm estimates. There is a fixed $A$ such that
\[
 B(q,a/A)\subseteq S_a(q)\subseteq B(q,Aa).
\]
Choose a sufficiently large fixed $K$, and increase $r_*$ to absorb $D$,
all lattice endpoint conventions, and the minimum admissible inner radius
$n_0(6)$ for the six-arm comparison.
When $t\ge4K^2r$, put
\[
 a=\lceil Kr\rceil+c_{\rm lat},\qquad
 b=\lfloor t/K\rfloor-c_{\rm lat},
\]
where the fixed integer $c_{\rm lat}$ absorbs the endpoint conventions of
the arm estimates.
Every original arm has a point inside $S_a(q)$ and reaches outside
$S_b(q)$. Retain the segment from its last visit to the inner boundary
before its first visit to the outer boundary, up to that outer visit.
Fixed lattice endpoint adjustments are absorbed by $c_{\rm lat}$.
This gives six disjoint arms in the lattice annulus, with
\[
 n_0(6)\le a<b\le t/K\le L,
 \qquad a/b\le C r/t;
\]
increasing $r_*$ once more ensures these inequalities despite rounding.

Their colors alternate: the two separated half-planes meet the outer boundary
of the common convex ball in two disjoint arcs.  Each group of three therefore occurs
consecutively in cyclic order.  Disjoint paths preserve the linear order
within a group under the annular truncation.  Its word is $BWB$, while
the other group's word is $WBW$, or conversely.  Reversing a group's
linear orientation does not change either word.  Their cyclic
concatenation is consequently alternating.

The critical six-arm gap holds in both models
\cite[Remark~4.6]{GPS}.  In the triangular regime,
\cite[Theorem~27]{Nolin} compares this alternating six-arm event under
$\nu$ to criticality at the radii $a,b$. The bound $b\le t/K\le L$
ensures that this comparison is within the characteristic-length cutoff.
Thus, in either model,
\[
 \Prob(E)^2\le C\alpha_{6,\mathrm{crit}}(a,b)
 \le C(a/b)^{2+\eta_6}
 \le C(r/t)^{2+\eta_6}.
\]
Taking square roots proves \eqref{eq:reflection}; the bounded-ratio
range follows from probability at most one.  If both original orders
are allowed, apply the argument separately and add the two bounds.
Increasing the fixed constants proves the padded versions as well.
\end{proof}

\begin{corollary}[Half-plane arms with unspecified order]\label{cor:generalreflection}
Fix an integer $j\ge1$.  In regime (ii), for every direction of $H$ and
$z\in\partial H$,
\[
 \Prob\bigl(A_j^{H,\mathrm{any}}(z;r,t)\bigr)
 \le C_j\sqrt{\alpha_{2j,\mathrm{crit}}(r,t)},
 \qquad r_*(j)\le r<t\le L_\varepsilon(p_*),
\]
where the $j$ site arms are disjoint and monochromatic but their colors
and order are unspecified.  The critical event on the right is alternating.
The constants have the same direction, phase, and padding uniformity as
Lemma~\ref{lem:reflection}.
\end{corollary}
\begin{proof}
For $0\le k\le j$, let $E_k$ be the event of $j$ disjoint arms with
$k$ open arms and $j-k$ closed arms.
Apply the site-inversion construction to $E_k$ and its color-reversed
copy. Then $\Prob(E_k)^2$ is bounded by the probability of $2j$
disjoint full-plane arms with exactly $j$ arms of each color, after
truncation to a common lattice annulus below the characteristic length.
Sum over the finitely many possible balanced color words.  For each word,
use \cite[Theorem~27]{Nolin} to compare with criticality and
\cite[Proposition~20]{Nolin} to compare the critical nonconstant word
with the alternating word.  Critical fixed-factor regularity absorbs
the radius changes.  Taking square roots and then summing over
$k=0,\ldots,j$ proves the result.
\end{proof}

Lemma~\ref{lem:reflection} gives \eqref{eq:halfarm} in both regimes with $\beta=1+\eta_6/2>1$, supplying the half-plane input in Theorem~\ref{thm:armcriterion}.

The pivotal scale and its boundary-arm compensation are recorded in \cite[Remark~36, equation~(7.17)]{Nolin}. We state the uniform interior estimates needed here.

\begin{lemma}[Bulk pivotal estimates]\label{lem:bulkmoments}
Fix a positive relative distance $\gamma$ from the sides of $Q$. Uniformly for relevant coordinates $x,y$ in this interior region,
\begin{align}
 c\alpha_4(R)\le\Prob(x\in\cP_f)&\le C\alpha_4(R),\label{eq:bulkone}\\
 \Prob(x,y\in\cP_f)&\le C\alpha_4(R)\alpha_4(|x-y|),\quad x\ne y.\label{eq:bulktwo}
\end{align}
The constants may depend on $\gamma$ and the previously fixed parameters. Consequently
\begin{equation}\label{eq:meancomparison}
 I(f)\asymp m(R),\qquad J(f)\asymp I(f).
\end{equation}
\end{lemma}
\begin{proof}
Pivotality at $x$ forces four alternating arms to radius $c_\gamma R$, for a sufficiently small fixed $c_\gamma>0$, giving the one-point upper bound. For the lower bound, require four arms from the coordinate to a square of radius $c_\gamma R$, with separated landings directed towards the four distinguished boundary arcs. Arm separation and comparison give probability at least $c\alpha_4(R)$. Extend them through four separated corridors to the respective sides of $Q$, open towards the left and right and closed towards the top and bottom. Uniform RSW and separated-arm extension bound the conditional probability of these extensions below by a positive constant, uniformly in the profile and scale. For mixed colors this gluing is justified by the separated-domain form of positive association, as in \cite[Lemma~13 and Theorem~11]{Nolin}; in the critical bond setting the corresponding arm-separation construction is reviewed in \cite[Section~2.2]{GPS}. The resulting outside configuration makes $x$ pivotal. Finite lattice endpoint adjustments are absorbed by finite energy and the fixed microscopic cutoff.

For the pair estimate put $d=|x-y|$. When $d$ is small compared with $R$, use two disjoint four-arm annuli of outer radius $cd$ about $x,y$, and a disjoint common four-arm annulus from radius $Cd$ to $c_\gamma R$. Independence and the full-plane arm comparison give
\[
 \Prob(x,y\in\cP_f)\le C\alpha_4(d)^2\alpha_4(d,R)
 \asymp C\alpha_4(d)\alpha_4(R).
\]
